\documentclass[12pt]{article}
\usepackage[utf8]{inputenc}
\usepackage[T1]{fontenc}
\usepackage{hyperref}
\title{Sufrai Language Resources for K-12 Curriculum Development}
\author{Compiled by Claude Sonnet 5.5 (claude-sonnet-5-5) for Andr\'as Kornai}
\date{}
\begin{document}
\maketitle

\section{Introduction}
Sufrai (ISO 639-3: suf; Glottocode tarp1241) appears in the search results under the name Tarpia, also called Tarfia or Kamdera. The sources describe it as an Austronesian language of the north coast of Papua province, Indonesia (Jayapura Regency). Most sources place it in the Sarmi branch of North New Guinea. Joshua Project instead links it to the Lakes Plain family, so the classification is disputed. Speaker counts are 630 (2006, Wikipedia) and 564 (Endangered Languages Project, data from 1978). Both sources label it endangered. Younger people are reported to be shifting to Indonesian. The ``Bookkeeping'' family label in the request does not match anything in these sources.

\section{Primary Educational Resources}
None found. The search turned up no primers, textbooks or school materials.

\section{Linguistic and Reference Materials}
\begin{itemize}
\item Wikipedia, Tarpia language: \url{https://en.wikipedia.org/wiki/Tarpia_language}
\item Endangered Languages Project entry: \url{https://endangeredlanguages.com/lang/2790}
\item Endangered Languages Project bibliography: \url{https://www.endangeredlanguages.com/lang/2790/bibliography?hl=de}
\item Verbix endangered-language page: \url{https://docs.verbix.com/EndangeredLanguages/Tarpia}
\end{itemize}

\section{Religious and Community Resources}
Joshua Project describes the Tarpia people of Tarfia village (Demta Subdistrict) and notes that no complete Bible translation exists: \url{https://joshuaproject.net/people_groups/15275}

\section{Digital and Audio Resources}
None found.

\section{Curriculum Development Context}
Very little exists. A curriculum would have to start from new documentation work with the community, and Indonesian would be the likely language of instruction in the meantime.
\end{document}